\documentclass[12pt,a4paper]{article}
\newcommand{\LabNumber}{6}
\newcommand{\LabTitle}{Funcții hash și semnături digitale:\\ RSA și ElGamal cu HAVAL-224,4}
\newcommand{\LabTitleShort}{Lucrarea de laborator nr. 6 — Criptografie \& Securitate Informațională}
% ---------------------------------------------------------------------------
% Common preamble for the CS laboratory reports (compile with XeLaTeX)
% Strunga Daniel-Ioan, FAF-242 — https://github.com/don09dandidi/CS
% ---------------------------------------------------------------------------
\usepackage{fontspec}
\setmainfont{Latin Modern Roman}
\setsansfont{Latin Modern Sans}
\setmonofont{DejaVu Sans Mono}[Scale=0.85]
\usepackage[english]{babel}
\babelprovide[import, main]{romanian}
\usepackage[a4paper, margin=2.5cm, headheight=15pt]{geometry}
\usepackage{amsmath, amssymb}
\usepackage{booktabs, array, longtable, tabularx}
\usepackage{xcolor}
\usepackage{listings}
\usepackage{fvextra}
\usepackage[most]{tcolorbox}
\usepackage{fancyhdr}
\usepackage{tikz}
\usetikzlibrary{arrows.meta, positioning, shapes.geometric, fit, calc}
\usepackage{enumitem}
\usepackage{caption}
\usepackage{float}
\usepackage[hidelinks]{hyperref}
\usepackage{url}

\definecolor{codebg}{RGB}{248,248,250}
\definecolor{codeframe}{RGB}{200,200,210}
\definecolor{kw}{RGB}{0,0,160}
\definecolor{cmt}{RGB}{0,120,0}
\definecolor{str}{RGB}{160,30,30}
\definecolor{termbg}{RGB}{30,32,40}
\definecolor{termfg}{RGB}{225,228,235}
\definecolor{termbar}{RGB}{60,63,75}
\definecolor{boxblue}{RGB}{230,238,250}

\lstdefinestyle{py}{
  language=Python, basicstyle=\ttfamily\small, backgroundcolor=\color{codebg},
  frame=single, rulecolor=\color{codeframe}, numbers=left, numberstyle=\tiny\color{gray},
  keywordstyle=\color{kw}\bfseries, commentstyle=\color{cmt}\itshape, stringstyle=\color{str},
  breaklines=true, showstringspaces=false, columns=fullflexible, keepspaces=true, tabsize=4,
  extendedchars=true, inputencoding=utf8
}
\lstdefinestyle{sh}{
  language=bash, basicstyle=\ttfamily\small, backgroundcolor=\color{codebg},
  frame=single, rulecolor=\color{codeframe}, numbers=left, numberstyle=\tiny\color{gray},
  keywordstyle=\color{kw}\bfseries, commentstyle=\color{cmt}\itshape, stringstyle=\color{str},
  breaklines=true, showstringspaces=false, columns=fullflexible, keepspaces=true
}
\lstdefinestyle{wl}{
  basicstyle=\ttfamily\small, backgroundcolor=\color{codebg},
  frame=single, rulecolor=\color{codeframe}, numbers=left, numberstyle=\tiny\color{gray},
  morekeywords={RandomPrime,RandomInteger,PowerMod,Mod,GCD,FromDigits,IntegerDigits,
    ToCharacterCode,FromCharacterCode,IntegerLength,BitLength,Module,Print,Hash,Length},
  keywordstyle=\color{kw}\bfseries, commentstyle=\color{cmt}\itshape,
  morecomment=[s]{(*}{*)}, breaklines=true, columns=fullflexible, keepspaces=true
}

% Console session box (as in the Lab 2 report)
\newtcolorbox{terminal}[1][Sesiune Terminal Consolă]{
  enhanced, breakable, colback=termbg, colframe=termbar, coltext=termfg,
  fonttitle=\small\ttfamily\bfseries, title={\textbullet\textbullet\textbullet\quad #1},
  left=4pt, right=4pt, top=2pt, bottom=2pt, boxrule=0.6pt, arc=2pt
}
\DefineVerbatimEnvironment{termout}{Verbatim}{fontsize=\scriptsize, breaklines=true, breakanywhere=true}

\newtcolorbox{infobox}[1]{
  enhanced, breakable, colback=boxblue, colframe=kw!60, fonttitle=\bfseries,
  title={#1}, boxrule=0.6pt, arc=2pt
}

\newcommand{\bignum}[1]{\par\nobreak\vspace{1pt}{\noindent\ttfamily\footnotesize\seqsplit{#1}\par}\vspace{5pt}}
\usepackage{seqsplit}

\pagestyle{fancy}
\fancyhf{}
\renewcommand{\headrulewidth}{0.4pt}
\fancyhead[L]{\footnotesize\textit{\LabTitleShort}}
\fancyhead[R]{\footnotesize strunga.daniel-ioan}
\fancyfoot[C]{\thepage}

\newcommand{\makereporttitle}{%
\begin{titlepage}
  \centering
  {\large\bfseries UNIVERSITATEA TEHNICĂ A MOLDOVEI\par}
  \vspace{2pt}
  {\normalsize Facultatea Calculatoare, Informatică și Microelectronică\par}
  {\normalsize Departamentul Ingineria Software și Automatică\par}
  \vspace{5cm}
  {\LARGE\bfseries LUCRAREA DE LABORATOR NR. \LabNumber\par}
  \vspace{8pt}
  {\Large\bfseries \LabTitle\par}
  \vspace{10pt}
  {\normalsize Disciplina: Criptografie și Securitate Informațională\par}
  \vfill
  \begin{tabular*}{\textwidth}{@{}l@{\extracolsep{\fill}}r@{}}
    \textbf{Autor:} & \textbf{Verificat de:} \\
    Daniel-Ioan Strunga & Lector / Supervizor \\
    Student gr. FAF-242 & Maia Zaica \\
  \end{tabular*}
  \vspace{3cm}

  {\normalsize Chișinău — 2026\par}
\end{titlepage}
}


\begin{document}
\makereporttitle
\tableofcontents
\newpage

\section{Introduction and Laboratory Objectives}
A \textbf{digital signature} gives a document the properties of a handwritten signature —
authenticity, integrity and non-repudiation — using public-key cryptography. Because public-key
operations are slow and work on numbers smaller than the modulus, the document itself is not
signed: a \textbf{cryptographic hash function} first compresses it into a short fixed-length
digest, and the digest is signed (\emph{hash-then-sign}).

The variant is determined by the student's number in the group list, $k = 22$:
\[
  i = (k \bmod 24) + 1 = (22 \bmod 24) + 1 = 23
\]
In \emph{both} lists of the assignment (Task 2 and Task 3) index 23 corresponds to
\textbf{HAVAL-224,4} — the HAVAL hash with a 224-bit digest and 4 passes.

Objectives:
\begin{enumerate}
  \item \textbf{Task 2} — generate RSA keys with $n \ge 3072$ bits, sign the message $m$ obtained
        in Laboratory Work No.~2 and validate the signature, using HAVAL-224,4.
  \item \textbf{Task 3} — sign and validate the same message with the \textbf{ElGamal} signature
        scheme, using the given 2048-bit prime $p$ and generator $g=2$, and HAVAL-224,4.
  \item Use the decimal representation of the hash value and comment every step of the algorithms.
\end{enumerate}

\begin{infobox}{Tooling note}
The guide suggests Wolfram Alpha / Mathematica. Wolfram's \texttt{Hash[]} function does not
implement HAVAL, so the digest was computed with a pure-Python implementation of HAVAL
(\texttt{lab6/haval.py}, checked against PHP's \texttt{hash("haval224,4", ...)}) and then
converted from hexadecimal to decimal. The modular arithmetic was performed with
\texttt{lab6/digital\_signatures.py} (Python integers, the same operations as \texttt{RandomPrime},
\texttt{RandomInteger}, \texttt{PowerMod}, \texttt{Mod}); the equivalent Wolfram Language code is
given in Appendix~\ref{app:wl}.
\end{infobox}

\section{The Message}
The message $m$ is the plaintext recovered in Laboratory Work No.~2 (Variant 22, a fragment of
D.~Kahn's \emph{The Codebreakers} about William F.~Friedman). It is stored in
\texttt{lab6/message\_lab2.txt} as UTF-8 text, three paragraphs separated by a newline, without a
trailing newline: \textbf{3555 bytes}.
\begin{infobox}{$m$ (beginning and end)}
\small\itshape While Herbert Yardley may be the best known cryptologist, uncontestably the greatest
is William Frederick Friedman. Unlike his contemporary, his eminence is due most emphatically to
what he did. [\dots] A cardinal article of faith with him was that a well-executed sales campaign
could put across almost anything.
\end{infobox}

\section{The HAVAL Hash Function}
\subsection{Description}
HAVAL was proposed by Y.~Zheng, J.~Pieprzyk and J.~Seberry (AUSCRYPT'92). It is a
Merkle–Damgård hash of the MD family, with two parameters: the number of passes (3, 4 or 5) and
the digest length (128, 160, 192, 224 or 256 bits).
\begin{itemize}
  \item \textbf{State:} eight 32-bit words, initialised with the first 256 bits of the fractional
        part of $\pi$:\\ {\small\texttt{243F6A88 85A308D3 13198A2E 03707344 A4093822 299F31D0 082EFA98 EC4E6C89}}.
  \item \textbf{Padding:} append byte \texttt{0x01}, then zeros up to 944 bits mod 1024, then
        16 bits encoding VERSION=1, PASS and FPTLEN, then the 64-bit message length
        (little-endian). The message is processed in \textbf{1024-bit blocks} of 32 words.
  \item \textbf{Compression:} each pass has 32 steps. Step $j$ of pass $r$ updates one register:
        \[ T = F_r\big(\phi_{P,r}(x_6,\dots,x_0)\big),\qquad
           x_7 \leftarrow (T \ggg 7) + (x_7 \ggg 11) + W_{\mathrm{ord}_r(j)} + K_{r,j}, \]
        where $F_1\dots F_5$ are highly non-linear Boolean functions of 7 variables, $\phi_{P,r}$
        is a fixed permutation of the inputs (depending on the number of passes $P$),
        $\mathrm{ord}_r$ a word order and $K_{r,j}$ further words of $\pi$ (none in pass 1). After
        all passes the result is added word-wise to the old state.
  \item \textbf{Tailoring (224 bits):} the 256-bit state is folded to 224 bits by adding 5- and
        4-bit fields of the last word $H_7$ to $H_0\dots H_6$.
\end{itemize}
HAVAL-224,4 therefore has $4\times32 = 128$ steps per block. Its security is lower than that of
SHA-2/SHA-3 (collisions for 3-pass HAVAL were found in 2004), so in practice it should not be
used for new signatures; here it is used as required by the variant.

\subsection{Implementation and Validation}
\begin{lstlisting}[style=py, caption={Compression function of HAVAL (\texttt{haval.py}, excerpt)}]
def _compress(state, block: bytes, passes: int):
    w = [int.from_bytes(block[4*i:4*i+4], "little") for i in range(32)]
    t = list(state)
    for p in range(passes):
        f, phi = _F[p], _PHI[passes][p]
        for s in range(32):
            x = [t[(j - s) % 8] for j in range(8)]        # x7..x0 rotate each step
            temp = f(*[x[j] for j in phi]) & MASK
            new = (_rotr(temp, 7) + _rotr(x[7], 11) + w[ORDER[p][s]]) & MASK
            if p > 0:
                new = (new + CONST[p][s]) & MASK
            t[(7 - s) % 8] = new
    return [(a + b) & MASK for a, b in zip(state, t)]
\end{lstlisting}
The implementation was compared with PHP's \texttt{hash()} for HAVAL-224 and HAVAL-256 with 3,
4 and 5 passes on 10 inputs of different lengths (including the block boundaries 117–129 bytes):
all 60 digests are identical. For example
\texttt{HAVAL-224,4("")} $=$ \texttt{3e56243275b3b81561750550e36fcd676ad2f5dd9e15f2e89e6ed78e}.

\subsection{Digest of the Message}
\noindent\textbf{Hexadecimal digest} $\mathrm{HAVAL\text{-}224,4}(m)$:
\bignum{699A7B65152794B3D652C28987E742E25767CAE4DBDDA1F7096E118A}
\noindent\textbf{Decimal value} $h$ (223 bits), used in both signatures:
\bignum{11121340838540799761581205197693366182088800469145110944605870625162}
\textbf{Avalanche effect:} replacing a single letter (the first ``Friedman'' $\to$ ``Fridman'')
gives the completely different digest
\bignum{0DDF80AB87E92B503527A9F2EDDBDBB60BC047752798F3BEEB908771}

\section{Task 2 — RSA Digital Signature}
\subsection{Algorithm}
\begin{enumerate}
  \item \textbf{Key generation:} random primes $p, q$ of 1536 bits each, $n = pq$ (3072 bits),
        $\varphi(n) = (p-1)(q-1)$, $e = 65537$ with $\gcd(e,\varphi(n)) = 1$,
        $d = e^{-1} \bmod \varphi(n)$. Public key $(e,n)$, private key $(d,n)$.
  \item \textbf{Signing} (with the private key): $h = H(m)$, \quad $s = h^{d} \bmod n$.
        The signed message is the pair $(m, s)$.
  \item \textbf{Verification} (with the public key): compute $h = H(m)$ from the received message
        and $h' = s^{e} \bmod n$; the signature is valid iff $h' = h$.
\end{enumerate}
Since $h < 2^{224} < n$, the digest is signed directly as a number (\emph{textbook} RSA signature;
real systems apply PKCS\#1 v1.5 or PSS padding first).

\subsection{Key Generation}
\noindent\textbf{Prime} $p$ (1536 bits):
\bignum{2345573859860826350398489684824278427182026391492908800894570498274797519824232807763141247868628402605420473404579171588281115130635120387833386452156601161716642207377721031116769710429782669114956776543624769674662747468843845205935233545127854676530730699828824187425293711003992170689742024165955844812931288141862494290608575286733958435086677555278225595459122725365223856052386298868427522638131235988658427753332331206307319117371641380563259206697755421}
\noindent\textbf{Prime} $q$ (1536 bits):
\bignum{1913764818546337516387054075315544084950321358420621482846311385627276986742964473954859886611119175202062895267680857123949048721354406649953185218869568370484349773164209582233860311052104717961544705682476222877659075331437336724012867829181541902901910526564605275138719630664973373482272292232709213876915434538320208552969330533461093260151583112702672811673492488319668415661714812439942941985717133152871292123631677064994722568591828228699418355227689019}
\noindent\textbf{Modulus} $n = p\cdot q$ (\textbf{3072 bits}):
\bignum{4488876732303586842898729734549412759818147223740736568339078506675704964313704934238321506438911857302513362707213089943861661966796556814208369476422886061679704530933030800165826174516323048264978226140813681740780928295495697374765163663655474348160820809858269672944888241642413860873213334623713122377377356269448894652615536372450083166527478028804586759937207968721671897338326463362765143882365130130534034859197639671604022297901262000894840443341787123061854171656175979315045837027916860215705864717960035314975093229917711364400367549737147877285331719124891989880480021934304019501365033066018187390499584543063715598261152206812083694467852692014258780240255033787515776544515133312545857197368073663626284471815148948740950105392137334950674967359575363726940237866743524867833144051439472968169267140453143783285286424110070494224393942118352669727459397038432317852950794815438062535227222963882616309421999}
\noindent\textbf{Euler's function} $\varphi(n)$:
\bignum{4488876732303586842898729734549412759818147223740736568339078506675704964313704934238321506438911857302513362707213089943861661966796556814208369476422886061679704530933030800165826174516323048264978226140813681740780928295495697374765163663655474348160820809858269672944888241642413860873213334623713122377377356269448894652615536372450083166527478028804586759937207968721671897338326463362765143882365130130534034859197639671604022297901262000894840443341787118802515493249012112529502076888094348083358114804429751574093209327843204797203085831736013397537753911641523317620451309704140167511837995279446516364330052342071735056330538856182062212580465615512776554139262481465692976263333203364444482887971494230985058078385686384727608436426593162936358568694516673880217557684040681289927323856387777729908599159554736650670072739217798780123282633747888045879090255508712440888942523513396376571757613701205054383977560}
\noindent\textbf{Public exponent} $e = 65537$; \textbf{private exponent} $d = e^{-1} \bmod \varphi(n)$:
\bignum{3404620225713938556845250147477724950063027664532718809259358459058688476145403835485058014870363179439645264813577689882654555920831842311424926706510099599736226454050203121654069024469885251996693054103505282180890147592721739313783847143282766428491232738691502672277821075534727951880995700507208266689233505471496959068885674755716256221044316193566833911778061194458826998504603407454918092176338915778849432698264142013769643657197888677822906662147950200868465715899852679822145654147680024416428609984951854089956072997834844146887617489938538809396358296625802211711853964195243836405556116870705830125299508854042140080336635716072474577730704859091712836666315518962039775563750301961280527197039825804348296106051807576264632689159660441431215578590694421602514215478288755128834360543349058815944683742374342656741341618449106977182172053583567619764590068672224381635818942220125939396865216659990534175570633}

\subsection{Signing and Verification}
\noindent\textbf{Signature} $s = h^{d} \bmod n$:
\bignum{4025922501518285721408071051742966348424968176118599527770000819092139797383596225837317623990142461152843519657441533062126248465668876377624787709280636657662141461097809723885232590576174807354836165510089558430648639473673406250041829025733695307628168893383187435407230987457597879592092749294536705833121293450017285263668911553638250133587947615176249708836334933412751016940003717367772300371036442095957895245739272118310006731272078102055197459396559990786627134362648447968896046891146897333529644877729677522036434552400803504483957381328032954461567637922962646881678129811816496409321469186998720482347812872467587770310911687920108218377948434906785537617488830602438855654447854490465392490429210118160226423601825475757796580778230974054486773902041792624141805496460139975626123395491685486072285626273176545769132677595213337504836434095702212872292404461439437486407742055727919817866744653935105079818863}
\noindent\textbf{Verification} $h' = s^{e} \bmod n$:
\bignum{11121340838540799761581205197693366182088800469145110944605870625162}
$h' = h$ $\Rightarrow$ the signature is \textbf{valid} (True). For the modified message
$m'$, $h' \ne H(m')$ $\Rightarrow$ the signature is rejected (check $=$ False).

\section{Task 3 — ElGamal Digital Signature}
\subsection{Algorithm}
Public parameters: the given 2048-bit prime $p$ and the generator $g = 2$.
\begin{enumerate}
  \item \textbf{Key generation:} private key $x$ random, $1 < x < p-1$; public key
        $y = g^{x} \bmod p$.
  \item \textbf{Signing:} choose a random $k$, $1<k<p-1$, with $\gcd(k, p-1) = 1$ (so that
        $k^{-1} \bmod (p-1)$ exists) and compute
        \[ r = g^{k} \bmod p, \qquad s = (h - x\,r)\,k^{-1} \bmod (p-1). \]
        The signature is the pair $(r,s)$.
  \item \textbf{Verification:} accept iff $0<r<p$ and
        \[ g^{h} \equiv y^{r}\, r^{s} \pmod p. \]
\end{enumerate}
\textbf{Correctness:} from the definition of $s$, $h \equiv x r + k s \pmod{p-1}$, therefore
$g^{h} = g^{xr}\,g^{ks} = (g^{x})^{r}(g^{k})^{s} = y^{r} r^{s} \pmod p$ (Fermat's little theorem
allows exponents to be reduced modulo $p-1$). The value $k$ must be secret and unique for every
signature — reusing it reveals the private key $x$.

\subsection{Values}
\noindent\textbf{Prime} $p$ (given, 2048 bits), $g = 2$:
\bignum{32317006071311007300153513477825163362488057133489075174588434139269806834136210002792056362640164685458556357935330816928829023080573472625273554742461245741026202527916572972862706300325263428213145766931414223654220941111348629991657478268034230553086349050635557712219187890332729569696129743856241741236237225197346402691855797767976823014625397933058015226858730761197532436467475855460715043896844940366130497697812854295958659597567051283852132784468522925504568272879113720098931873959143374175837826000278034973198552060607533234122603254684088120031105907484281003994966956119696956248629032338072839127039}
\noindent\textbf{Private key} $x$:
\bignum{10875856765490333315570046701595855041281648793126689674587727305962167509626780176716572916399412072063036222923455297292064259062373111516161442474582544100734476024049313471473577760215629997839867376658698294882040363828106357948964710350157832310705623534043431538516842463454201829997630882467256049509778005963602485262590757137426443164091705189097790715922680413139310599272870357986449488849223426616350056766691007286901496461997068004948648171346684739891080760177692746844411007464213614348571261019245672950133918032627330668243612474228745209247313220142413789027555770677685682272036639249191386637351}
\noindent\textbf{Public key} $y = g^{x} \bmod p$:
\bignum{3755878000369411850106859280579880329889281561075821269562508368359796020389619863003747982010543681986036442539519218506278551260763853019636130596455403975202936646925354436583695344299748455736506471894685572299324597804020528920064282996828787956228926801992860922096253488423420838199874848228562156297722699530628719555094806013145576509420099967439840993714270805346704450584333138224679373330256969193712683133971356939128367025114612262167124835579328295128125561153274469154088285559338881504906588018728419669466625682616952068414216814636408703244783119192660425817403243355802946318646436581015185456801}
\noindent\textbf{Random} $k$ ($\gcd(k,p-1) = 1$):
\bignum{104710444496975478300862294525129905276030212534197903365870526408454444756390800396486985541030739117359426874676064268544759231389442339408096214263470485542013201828672921042086323903643025685198666968114522130302751261308511707364919525961448565695945811605365533906609543692313002439829556526913814275155972302119464551650434016907808254730385170052292802151623292700622480704621580503934385604749872361644820635953960841396459214232865315671921027231041925646523949285910391096111007750717494104918291235198879677509963372816920666338474821225497148856874869775789054418379251504412307490047050141022744332973}
\noindent$k^{-1} \bmod (p-1)$:
\bignum{8383502745810764520165380568982957255879357487883062011257920554843409374762876061063248881232465700975184666431994991435487512967327826344070229957306705927601708673144174526725604875217118362942172611397052044295119192897067561833977215739964408662029675665752873756183858852751277403604877076848114939673187989091124722372973313321437255820380272580312479355042605688030786207347447904479009452606092025253931372593043928127225888614932593068561567768487261022665301703630626294288605258581116397717269276187037386023922482229581644189273981699536460998588477592644012293032748715754997381210725181170263109448451}
\noindent\textbf{Signature} $r = g^{k} \bmod p$:
\bignum{23034125902480498625022098289671025249083559319889468590753155582003319406962591889693518105543938956563611824169438537258782941211361763841810998969462854230678981894170081896834785211423219029830497747349015669425940134588883443618298317713394441988171200916251558242638891527074426902681065186871387073003402665286581471258005219784733289948934308393524480372487637229214387621019952008635740927193724968804491407412049066424244858082457555494332005036977145190957216122022539389122406587835617915443993398710428911174302576103312738165925538250281226242331689563829528330543949331931112399057642345579443228286446}
\noindent\textbf{Signature} $s = (h - x r)\,k^{-1} \bmod (p-1)$:
\bignum{9590267148580090127401284866732798650884341560560935929011696907286271272313894286838058066648430107109351631033028076788116742191070064830432524537146220155194715925045036583060503418833525871378967272380214797239442262665478776936745836607050611130453605268230273809784550724373743445313128259110381247071823162278160145383077355565566909449957214096699747517699727740459883832339145844894034441887238766511176446265588033022531041618361122494652144150697689379083283300907961114109882931643060979360041502188933541479196248806535706554605790829235443330661566561260256681200114220076322032560798295491859713419798}

\subsection{Verification}
\noindent Left side $g^{h} \bmod p$:
\bignum{5873012416190521668247857640938187909096522834311699887536701935616727645502198487623053713751149264027029046987116921308931307116273823003563716247664919500314419575129699123820621910792002283092085947467659397268543006954146663422117698441316738554761291852726707453665192322977727299715933640093462093403712064682526671811762128782289783983358758743259690742861865401603616689822200457727403834812671962751631014859761060063883708652098934088154491484006203995130792309891699774547243385863678240551525921777980279115073137038273978993735352810205345554634099486215967668661113754353599642779821624881917011341416}
\noindent Right side $y^{r} r^{s} \bmod p$:
\bignum{5873012416190521668247857640938187909096522834311699887536701935616727645502198487623053713751149264027029046987116921308931307116273823003563716247664919500314419575129699123820621910792002283092085947467659397268543006954146663422117698441316738554761291852726707453665192322977727299715933640093462093403712064682526671811762128782289783983358758743259690742861865401603616689822200457727403834812671962751631014859761060063883708652098934088154491484006203995130792309891699774547243385863678240551525921777980279115073137038273978993735352810205345554634099486215967668661113754353599642779821624881917011341416}
The two sides are equal $\Rightarrow$ the signature is \textbf{valid} (True). For the
modified message $m'$, $g^{H(m')} \ne y^{r} r^{s}$ and the signature is rejected
(check $=$ False).

\section{Source Code and Repository}
\begin{center}
\url{https://github.com/don09dandidi/CS} \quad (directory \texttt{lab6/})
\end{center}
\begin{lstlisting}[style=py, caption={Signing and verification (excerpt from \texttt{digital\_signatures.py})}]
h = int(haval224_4(message_bytes), 16)          # decimal value of the digest

# RSA (n = 3072 bits)
s = pow(h, d, n)                                # sign
valid = pow(s, e, n) == h                       # verify

# ElGamal (p = 2048 bits, g = 2)
r = pow(G, k, P)
s = ((h - x * r) * pow(k, -1, P - 1)) % (P - 1)  # sign: (r, s)
valid = pow(G, h, P) == (pow(y, r, P) * pow(r, s, P)) % P
\end{lstlisting}

\subsection{Build and Execution Instructions}
\begin{lstlisting}[style=sh]
git clone https://github.com/don09dandidi/CS.git
cd CS/lab6
python3 haval.py < message_lab2.txt        # HAVAL-224,4 digest of the message
python3 digital_signatures.py --save       # keys, signatures, verification -> results.json
\end{lstlisting}

\section{Conclusion}
\begin{enumerate}
  \item For $k = 22$ the hash function is HAVAL-224,4; it was implemented from its specification,
        validated against an independent implementation, and produced the 224-bit digest of the
        Lab~2 message, used in decimal form.
  \item \textbf{RSA signature:} with a 3072-bit modulus the digest was signed with the private
        exponent and verified with the public one ($s^{e} \bmod n = h$).
  \item \textbf{ElGamal signature:} with the given 2048-bit prime and $g=2$ the signature
        $(r,s)$ was generated with a random $k$ coprime to $p-1$ and verified through
        $g^{h} \equiv y^{r} r^{s} \pmod p$.
  \item In both schemes a single-letter change of the message produces a completely different
        digest and the signature is rejected, demonstrating integrity protection. RSA signatures
        are deterministic and short to verify ($e$ small); ElGamal signatures are randomised and
        twice as long, and their security depends critically on never reusing $k$.
\end{enumerate}

\appendix
\section{Wolfram Language Code}\label{app:wl}
\lstinputlisting[style=wl]{wolfram/lab6.wl}

\begin{thebibliography}{9}
\bibitem{haval} Y.\ Zheng, J.\ Pieprzyk, J.\ Seberry, \emph{HAVAL — A One-Way Hashing Algorithm
  with Variable Length of Output}, AUSCRYPT'92, LNCS 718, 1993.
\bibitem{rsa} R.\ Rivest, A.\ Shamir, L.\ Adleman, \emph{A Method for Obtaining Digital Signatures
  and Public-Key Cryptosystems}, Communications of the ACM, 21(2), 1978.
\bibitem{elg} T.\ ElGamal, \emph{A Public Key Cryptosystem and a Signature Scheme Based on Discrete
  Logarithms}, IEEE Trans.\ Inf.\ Theory, 31(4), 1985.
\bibitem{kahn} David Kahn, \emph{The Codebreakers: The Story of Secret Writing}, Scribner, 1996.
\bibitem{stallings} William Stallings, \emph{Cryptography and Network Security: Principles and
  Practice}, 7th Edition, Pearson, 2017.
\end{thebibliography}
\end{document}
